Tato část popisuje praktické postupy pro živé ukázky a interaktivní simulace v hodině se zaměřením na světy typu Minecraft Education a další nástroje, které snadno nasadíte pro názorné předvedení principů agentního chování, učení a diagnostiky chyb.

Krátké shrnutí přínosu
\begin{itemize}
  \item Interaktivní simulace umožňují studentům pozorovat chování „agenta“ v kontrolovaném prostředí a zkoušet dopady jednoduchých změn pravidel či dat (general background knowledge).
  \item Minecraft Education nabízí hotové „světy“ a aktivity, které lze využít k vysvětlení principů strojového učení a k realizaci krátké „Hour of Code“ lekce dostupné zdarma; zbytek obsahu je součástí školní licence Microsoft 365 A3/A5 \cite{kb_ai_101_tipu_pdf_04e4a115_page_15_2}.
\end{itemize}

Praktické scénáře (konkrétní příklady)
\begin{enumerate}
  \item Demonstrace principů agentního chování v Minecraft Education: vyberte předpřipravený svět s aktivitou zaměřenou na AI nebo Hour of Code, spusťte krátkou ukázku (10--20 min), nechte studenty komentovat, proč agent reagoval tak jak reagoval, a diskutujte omezení přístupu. Doporučený start: stáhnout Minecraft Education a vyzkoušet jednu z volně dostupných lekcí Hour of Code \cite{kb_ai_101_tipu_pdf_04e4a115_page_15_2}.
  \item Diagnostika chyb a adaptivní procvičování: využijte nástroj typu „Pokrok v matematice“ jako ukázku, jak systém generuje příklady, sbírá postupy studentů a poskytuje přehledy silných/ slabých stránek třídy; při ukázce vysvětlete, proč je nahrání postupu důležité pro formativní hodnocení \cite{kb_ai_101_tipu_pdf_04e4a115_page_8_1}.
  \item Rychlá tvorba vizuálů pro diskusi: během debriefu si v reálném čase nechte vygenerovat ilustraci nebo schéma (např. pomocí Designer) k probíranému tématu a porovnejte různé stylistické varianty; to pomáhá udržet pozornost a ukazuje limity automatické vizualizace \cite{kb_ai_101_tipu_pdf_04e4a115_page_51_1}.
\end{enumerate}

Před hodinou — kontrolní seznam přípravy
\begin{itemize}
  \item Jasně definujte cíle sezení (co mají studenti pochopit/umět pozorovat). (general background knowledge)
  \item Ověřte technické prostředí: instalace Minecraft Education (nebo alternativy), přístup do potřebných webových služeb, záložní plán pro případ výpadku sítě. Adresa pro Minecraft Education a volné lekce je uvedena v dostupných materiálech produktu \cite{kb_ai_101_tipu_pdf_04e4a115_page_15_2}.
  \item Připravte krátký instrukční list pro studenty s pravidly použití nástrojů a jednou reflexní otázkou (např. „Co by změnila drobná úprava pravidel chování agenta?“). (general background knowledge)
  \item Rozhodněte, zda budou studenti pracovat individuálně, ve dvojicích nebo v malých skupinách; připravte časovou osu aktivity (celkem 30--60 min doporučeno).
\end{itemize}

Průběh ukázky — doporučený časový rámec (45--60 min)
\begin{enumerate}
  \item 5--10 min: úvod a stanovení učebního cíle.
  \item 15--25 min: live demo v Minecraft Education nebo jiném simulátoru (včetně krátkých experimentů s parametry).
  \item 10--15 min: studentské mini-aktivity (např. navrhnout změnu pravidla a předpovědět dopad).
  \item 10--15 min: debrief, diskuse, porovnání pozorování s předpověďmi, závěrečná reflexe.
\end{enumerate}

Facilitace a pedagogické tipy
\begin{itemize}
  \item Moderujte diskusi tak, aby studenti reflektovali nejen „co se stalo“, ale proč se to stalo a jaké jsou limity daného modelu/světa (general background knowledge).
  \item Při použití adaptivních cvičení (např. nástrojů typu Pokrok v matematice) zdůrazněte, že nahraný postup je klíčový pro pedagogickou diagnostiku; ukažte příklady, jak systém identifikuje chybu a navrhuje oblast pro další procvičování \cite{kb_ai_101_tipu_pdf_04e4a115_page_8_1}.
  \item Používejte generované vizuály (Designer apod.) jako doplněk, nikoli jako jediný zdroj vysvětlení; ukažte studentům, jak různé prompty mění výsledek \cite{kb_ai_101_tipu_pdf_04e4a115_page_51_1}.
\end{itemize}

Zhodnocení a reflexe
\begin{itemize}
  \item Na konci sezení požádejte studenty o krátké shrnutí (1--2 věty) o tom, co se naučili a co by změnili v experimentu (může sloužit jako rychlý formativní checkpoint). (general background knowledge)
  \item Zaznamenejte technické a etické problémy pro kvartální revizi pilotu (co fungovalo, kde vznikly nejasnosti, potřeba jiných licencí nebo DPIA pokud sbíráte citlivá data). Použití Minecraft Education a dostupnost Hour of Code materiálů jsou vhodné pro první nízkoprahové piloty \cite{kb_ai_101_tipu_pdf_04e4a115_page_15_2}.
\end{itemize}

Krátké doporučení pro další kroky
\begin{itemize}
  \item Pro první pilot zvolte jednu ukázkovou lekci z Minecraft Education (Hour of Code) a proveďte ji včetně reflexe; výsledky pilotu použijte pro úpravu sylabu a instrukcí pro studenty \cite{kb_ai_101_tipu_pdf_04e4a115_page_15_2}.
  \item Pokud budete demonstrovat adaptivní procvičování (diagnostiku chyb), předveďte i způsob, jak učitel finalizuje automatické návrhy (uvedeno u „Pokrok v matematice“) \cite{kb_ai_101_tipu_pdf_04e4a115_page_8_1}.
\end{itemize}

Poznámka: část praktických a organizačních kroků je uvedena jako obecné pedagogické doporučení; konkrétní právní a GDPR aspekty je vhodné řešit podle kapitol věnovaných ochraně osobních údajů a transparentnosti v této příručce (viz příslušné šablony a checklisty).